\chapter{Στατιστικοί Έλεγχοι Σύγκρισης Μοντέλων}
\label{ch:dm}

Η αξιολόγηση βάσει μέσης απόλυτης τιμής σφάλματος (MAE) δεν διακρίνει εάν
οι διαφορές επίδοσης μεταξύ δύο μοντέλων είναι στατιστικά σημαντικές ή
απλώς τυχαία αποτελέσματα συγκεκριμένης περιόδου δοκιμής. Στο κεφάλαιο αυτό
εφαρμόζεται ο έλεγχος Diebold-Mariano με τη διόρθωση Harvey-Leybourne-Newbold
για πλήρη ζευγαρωτή σύγκριση 18 μοντέλων ανά εργασία.

% ─────────────────────────────────────────────────────────────────────────────
\section{Ο Έλεγχος Diebold-Mariano}
\label{sec:dm_theory}
% ─────────────────────────────────────────────────────────────────────────────

Έστω $e_{i,t}$ και $e_{j,t}$ τα σφάλματα πρόβλεψης δύο μοντέλων $i$ και $j$
σε χρόνο $t$. Ορίζουμε τη διαφορά απωλειών:
\begin{equation}
    d_t = L(e_{i,t}) - L(e_{j,t}), \quad L(e) = |e|
    \label{eq:dm_loss}
\end{equation}
Ο έλεγχος DM~\cite{diebold1995} εξετάζει την υπόθεση:
\begin{equation}
    H_0: \mathbb{E}[d_t] = 0 \quad \text{(ισοδύναμη πρόβλεψη)}
\end{equation}
Η στατιστική δοκιμής είναι:
\begin{equation}
    \text{DM} = \frac{\bar{d}}{\sqrt{\hat{V}(\bar{d}) / T}}
    \label{eq:dm_stat}
\end{equation}
όπου $\bar{d} = T^{-1} \sum_t d_t$ και $\hat{V}(\bar{d})$ εκτιμάται μέσω
long-run variance (Newey-West HAC estimator με $h-1$ lags για $h$-step-ahead
πρόβλεψη).

\textbf{Πρόσημο σύμβασης:} DM$(i,j) > 0$ σημαίνει ότι το μοντέλο $i$ έχει
χαμηλότερα σφάλματα (καλύτερο) από το $j$· DM$(i,j) < 0$ σημαίνει ότι το
$j$ είναι καλύτερο.

% ─────────────────────────────────────────────────────────────────────────────
\section{Διόρθωση Harvey-Leybourne-Newbold}
\label{sec:hln}
% ─────────────────────────────────────────────────────────────────────────────

Ο κλασικός DM έλεγχος υπολογίζει μια $t$-στατιστική που ακολουθεί κανονική
κατανομή για μεγάλα $T$. Σε μικρά δείγματα, υπάρχει τάση απόρριψης $H_0$
συχνότερα από το ονομαστικό επίπεδο σημαντικότητας. Η διόρθωση
Harvey-Leybourne-Newbold (HLN)~\cite{harvey1997} προσαρμόζει τη στατιστική:
\begin{equation}
    \text{DM}^* = \left[\frac{T + 1 - 2h + T^{-1} h(h-1)}{T}\right]^{1/2} \cdot \text{DM}
    \label{eq:hln}
\end{equation}
και συγκρίνει το αποτέλεσμα με $t_{T-1}$ κατανομή. Στην παρούσα εφαρμογή
($T = 2160$, $h = 1$ για CL ή $h = 24$ για MIMO), η HLN διόρθωση είναι μικρή
αλλά εφαρμόζεται για ακρίβεια.

% ─────────────────────────────────────────────────────────────────────────────
\section{Πειραματική Διάταξη}
\label{sec:dm_setup}
% ─────────────────────────────────────────────────────────────────────────────

\begin{table}[htbp]
\centering
\caption{Μοντέλα που συμμετέχουν στους DM ελέγχους (Q1 2026, $T=2160$ ώρες).}
\label{tab:dm_models}
\begin{tabular}{lll}
\toprule
\textbf{Κωδικός} & \textbf{Περιγραφή} & \textbf{Στρατηγική} \\
\midrule
MR-TF-Ens    & Ensemble TF Best3 (MR)  & Walk-Forward TF \\
MR-TF-LGBM  & LGBM Teacher-Forcing (MR) & Walk-Forward TF \\
MR-TF-XGB   & XGB Teacher-Forcing (MR) & Walk-Forward TF \\
MR-TF-RF    & RF Teacher-Forcing (MR) & Walk-Forward TF \\
MR-Rec-XGB-SS & XGB Recursive + SS (MR) & Walk-Forward Rec \\
CL-Ens      & Ensemble TF Best3 (static) & Teacher-Forcing \\
CL-LGBM     & LGBM static CL & Teacher-Forcing \\
CL-XGB      & XGB static CL & Teacher-Forcing \\
OL-LGBM-DO  & LGBM Daily Optuna (static) & Recursive \\
OL-XGB-SS-D & XGB SS Dense (static) & Recursive + SS \\
MIMO-Ens    & Ensemble MIMO Best3 (static) & MIMO H=24 \\
MIMO-LGBM   & LGBM Direct Dense (static) & MIMO H=24 \\
MIMO-XGB    & XGB MIMO Dense (static) & MIMO H=24 \\
MIMO-RF     & RF two-stage Dense (static) & MIMO H=24 \\
MR-MIMO-Ens & Ensemble MIMO Best3 (MR) & Walk-Forward MIMO \\
MR-MIMO-LGBM & LGBM Direct Dense (MR) & Walk-Forward MIMO \\
Naive-1     & Naive ημέρας (persistence) & Baseline \\
Seasonal    & Εποχικό Profile & Baseline \\
\bottomrule
\end{tabular}
\end{table}

Ο πίνακας αποτελείται από $18 \times 17 / 2 = 153$ ζεύγη. Ο έλεγχος εφαρμόζεται
αμφίπλευρα (two-sided) σε επίπεδο σημαντικότητας $\alpha = 0,05$ και $\alpha = 0,10$.
Τα $p$-values υπολογίζονται από την $t_{T-1}$ κατανομή με την HLN~\cref{eq:hln}.

% ─────────────────────────────────────────────────────────────────────────────
\section{Αποτελέσματα — Τιμή DAM}
\label{sec:dm_price}
% ─────────────────────────────────────────────────────────────────────────────

\begin{figure}[htbp]
    \centering
    \includegraphics[width=\textwidth]{ch5_dm_test/dm_price_heatmap.png}
    \caption{Χάρτης θερμότητας DM ελέγχου για τιμή DAM (Q1 2026). Κάθε κελί
    $(i,j)$ χρωματίζεται σκούρο μπλε αν το μοντέλο $i$ είναι στατιστικά
    σημαντικά καλύτερο από το $j$ ($p<0{,}05$), ανοιχτό μπλε για $p<0{,}10$,
    λευκό για μη σημαντική διαφορά.}
    \label{fig:dm_price_heatmap}
\end{figure}

\begin{figure}[htbp]
    \centering
    \includegraphics[width=0.75\textwidth]{ch5_dm_test/dm_price_strategy.png}
    \caption{Αριθμός στατιστικά σημαντικών «νικών» ($p<0{,}05$) ανά μοντέλο
    και στρατηγική για την τιμή DAM.}
    \label{fig:dm_price_strategy}
\end{figure}

Τα αποτελέσματα για την τιμή DAM (\cref{fig:dm_price_heatmap,fig:dm_price_strategy})
αναδεικνύουν ξεκάθαρη ιεράρχηση:

\begin{enumerate}
    \item \textbf{MR-TF-Ens} (Ensemble Walk-Forward TF): \textbf{νικά στατιστικά
    σημαντικά όλα τα υπόλοιπα 17 μοντέλα} ($p < 0{,}05$), επιβεβαιώνοντας ότι
    η Monthly Retraining είναι στατιστικά ανώτερη στρατηγική για Q1 2026.

    \item \textbf{Ομάδα Walk-Forward TF} (MR-TF-LGBM, MR-TF-XGB): Νικούν
    15--17 μοντέλα, με μικρές εσωτερικές διαφορές μεταξύ τους (μη σημαντικές).

    \item \textbf{Static CL} (CL-Ens, CL-LGBM): Νικούν 10--12 μοντέλα, υπερέχοντας
    των Recursive και MIMO static αλλά υποχωρώντας έναντι Walk-Forward.

    \item \textbf{Recursive/OL static}: 6--9 νίκες· ανταγωνιστικά με static CL
    σε μεμονωμένες συγκρίσεις.

    \item \textbf{MIMO static}: 3--6 νίκες· χαμηλότερη κατάταξη, κυρίως πάνω
    από baselines.

    \item \textbf{Naive-1, Seasonal Profile}: 0--2 νίκες· οριακά καλύτερα μόνο
    έναντι SVR που κατέρρευσε σε Q1 2026.
\end{enumerate}

\noindent\textit{Σημαντική παρατήρηση:} Οι εσωτερικές διαφορές εντός της ομάδας
Walk-Forward TF (LGBM vs XGB vs RF) δεν είναι στατιστικά σημαντικές, υποδηλώνοντας
ότι όλοι οι αλγόριθμοι ωφελούνται εξίσου από τη μηνιαία επανεκπαίδευση. Αντίστοιχα,
εντός της ομάδας static CL, LGBM και XGB δεν διαφέρουν σημαντικά.

% ─────────────────────────────────────────────────────────────────────────────
\section{Αποτελέσματα — Ηλεκτρικό Φορτίο}
\label{sec:dm_load}
% ─────────────────────────────────────────────────────────────────────────────

\begin{figure}[htbp]
    \centering
    \includegraphics[width=\textwidth]{ch5_dm_test/dm_load_heatmap.png}
    \caption{Χάρτης θερμότητας DM ελέγχου για ηλεκτρικό φορτίο (Q1 2026).}
    \label{fig:dm_load_heatmap}
\end{figure}

\begin{figure}[htbp]
    \centering
    \includegraphics[width=0.75\textwidth]{ch5_dm_test/dm_load_strategy.png}
    \caption{Αριθμός στατιστικά σημαντικών «νικών» ανά μοντέλο για το φορτίο.}
    \label{fig:dm_load_strategy}
\end{figure}

Για το ηλεκτρικό φορτίο (\cref{fig:dm_load_heatmap,fig:dm_load_strategy}),
η εικόνα διαφέρει ουσιαστικά:

\begin{enumerate}
    \item \textbf{CL-Ens} (static Ensemble Best3): Κορυφαίο μοντέλο με \textbf{16 νίκες}
    (έναντι 17 του MR-TF για τιμή). Η μηνιαία επανεκπαίδευση δεν αποδίδει
    τόσο εντυπωσιακά πλεονέκτημα για το φορτίο.

    \item \textbf{MR-TF-Ens}: 14--15 νίκες· στατιστικά ισοδύναμο με CL-Ens
    ($p = 0{,}12$ για ζεύγος CL-Ens vs MR-TF-Ens). Η Walk-Forward MR βελτιώνει
    αλλά η διαφορά δεν αγγίζει συχνά το κατώφλι $p < 0{,}05$.

    \item \textbf{MR-Rec-XGB-SS}: 13--14 νίκες, σχεδόν ισοδύναμο με MR-TF-Ens
    για φορτίο. Αυτό επιβεβαιώνει το εύρημα ότι το αναδρομικό σφάλμα για
    φορτίο είναι μικρό λόγω ομαλής δυναμικής.

    \item \textbf{MIMO static}: Σαφώς χαμηλότερη κατάταξη (4--7 νίκες)·
    η MIMO αδυνατεί να ανταγωνιστεί TF/CL για φορτίο λόγω μη πρόσβασης στα
    πραγματικά πρόσφατα δεδομένα.
\end{enumerate}

\medskip
\noindent\textbf{Βασικό συμπέρασμα σύγκρισης φορτίου vs τιμής:} Για \textit{τιμή},
η Walk-Forward MR είναι αδιαμφισβήτητα η καλύτερη στρατηγική (ξεπερνά τα πάντα
με $p<0{,}05$). Για \textit{φορτίο}, η CL και η Walk-Forward MR είναι στατιστικά
αδιάκριτες· η ομαλή δυναμική του φορτίου σημαίνει ότι η μηνιαία επανεκπαίδευση
δεν προσφέρει τόσο σαφές στατιστικό πλεονέκτημα.

% ─────────────────────────────────────────────────────────────────────────────
\section{Συμπεράσματα Κεφαλαίου}
\label{sec:dm_conclusions}
% ─────────────────────────────────────────────────────────────────────────────

Ο DM test με HLN διόρθωση επιβεβαιώνει, πέραν της απλής σύγκρισης MAE, τα
ακόλουθα στατιστικά σημαντικά ευρήματα:

\begin{itemize}
    \item Η Walk-Forward Monthly Retraining (TF) είναι στατιστικά ανώτερη
    όλων των εναλλακτικών για \textbf{τιμή} ($\alpha = 0{,}05$).
    \item Για \textbf{φορτίο}, η Walk-Forward MR και το static CL Ensemble
    είναι στατιστικά ισοδύναμα · επιλέγεται το απλούστερο (static CL)
    αν η επιχειρησιακή πολυπλοκότητα είναι παράγοντας.
    \item Εντός κάθε στρατηγικής, οι διαφορές LGBM/XGB/RF δεν είναι
    συστηματικά σημαντικές· το ensemble αποδεικνύεται συνεπώς η πλέον ασφαλής
    επιλογή.
    \item MIMO στρατηγικές είναι σαφώς κατώτερες CL/TF για αμφότερες τις
    εργασίες, στατιστικά επιβεβαιωμένα.
\end{itemize}
